\clearpage\SourceLeaf{517}
\renewcommand{\BookRunningTitle}{Hilfstafel zur Lösung der unbestimmten Gleichung}
\pdfbookmark[0]{Hilfstafel zur Lösung der unbestimmten Gleichung}{hulfstafel-ausschliessung}
\begin{center}
\textbf{HILFSTAFEL}\\[.7em]
ZUR LÖSUNG DER UNBESTIMMTEN GLEICHUNG
\[
A=fxx+gyy
\]
MIT DEM AUSSCHLUSSVERFAHREN.
\end{center}
\BookRule{.18\textwidth}
Es wird vorausgesetzt, dass als Ausschlusszahl eine Primzahl $p$ gewählt wurde, durch die keine der Zahlen $A$, $f$, $g$ teilbar ist. Die Tafel beschränkt sich außerdem auf die beiden Fälle, in denen der Wert des Ausdrucks $\frac{A}{f}(\mathrm{mod}.p)$ entweder ein bestimmter quadratischer Rest (stets $1$) oder ein bestimmter quadratischer Nichtrest modulo $p$ ist. Schließlich hat man sich damit begnügt, die Tafel nur für den Fall einzurichten, dass $fg$ ein quadratischer Rest von $p$ ist, und den entgegengesetzten Fall ganz übergangen. Der sechste Abschnitt der \textit{Disquisitiones Arithmeticae} gibt ausreichend Anleitung, wie sich das, was die Tafel nicht unmittelbar enthält, leicht aus ihr ableiten lässt.
\par\medskip\textit{Beispiele.} Die zu lösende Gleichung sei $21680143=xx+78yy$
\par\medskip\noindent 1) Ausschlusszahl $=5$\hfill $fg=\mathrm{N},\quad \frac{A}{f}\equiv3\equiv2\pdot2^2$
\par\noindent Aus der Tafel $1,4$ für $fg=R$ und daher $0,2,3$ für $fg=\mathrm{N}$,
und für den obigen Fall $0,1,4$ oder ausgeschlossen $5n\pm2$
\par\medskip\noindent 2) Ausschlusszahl $=7$\hfill $fg=R,\quad \frac{A}{f}\equiv2\equiv1,3^2$
\par\noindent Aus der Tafel $0,1,2,5,6$ Für den vorliegenden Fall $0,3,6,1,4$ und ausgeschlossen $7n\pm2$
\par\medskip\noindent 3) Ausschlusszahl $=11$\hfill $fg=R,\quad \frac{A}{f}\equiv1$
\par\noindent Man erhält also $0,1,3,5,6,8,10$; ausgeschlossen $11n\pm2,4$
\par\medskip\noindent 4) Ausschlusszahl $=17$\hfill $fg=\mathrm{N},\quad \frac{A}{f}\equiv9\equiv1\pdot3^1$
\[
\begin{gathered}
0\pdot1\pdot3\pdot4\pdot6\\
0\pdot3\pdot8\pdot5\pdot1\\
0,2,3,4,6,7\\
\text{Ausschlusszahl}\quad\pm1,5,8
\end{gathered}
\]
